\begin{proof}[Proof of \cref{obj:lem:causal-nonempty}]
Fix an admissible tuple \(v\). Thus
\[
1<p\le2,\qquad
0<\alpha,\beta\le1,\qquad
\frac14\le\gamma\le1.
\]
We establish the witness construction and distance bounds first, followed by the identification and completion properties.

\begin{enumerate}
\item Define the separated witness primitives by
\[
e_{\mathrm{sep}}(x)=\frac12+\frac{\sin(2\pi x)}{10},
\qquad
m_{0,\mathrm{sep}}(x)=\frac{\cos(2\pi x)}8,
\qquad
\tau_{\mathrm{sep}}(x)=\frac{\cos(2\pi x)}8,
\qquad x\in[0,1].
\]
For \(y_\ast\in\mathbb R\), write \(\delta_{y_\ast}\) for the unit point mass, defined by
\[
\delta_{y_\ast}(\mathsf E)=\mathbf 1_{\mathsf E}(y_\ast)
\qquad
\bigl(\mathsf E\in\mathcal B(\mathbb R)\bigr),
\]
and use the same notation for point masses on the treatment space \(\{0,1\}\).
The witness law is
\[
\begin{aligned}
P_{\mathrm{sep}}(dx,da,dy)
=\lambda(dx)\bigl[
 &(1-e_{\mathrm{sep}}(x))\,
 \delta_0(da)\delta_{m_{0,\mathrm{sep}}(x)}(dy)\\
 &+e_{\mathrm{sep}}(x)\,
 \delta_1(da)\delta_{m_{0,\mathrm{sep}}(x)+\tau_{\mathrm{sep}}(x)}(dy)
\bigr].
\end{aligned}
\]
The two weights are nonnegative and sum to one, so this construction is a probability law with covariate marginal \(\lambda\). Its arm kernels and means satisfy
\[
Q_{0,P_{\mathrm{sep}}}(dy\mid x)
=\delta_{m_{0,\mathrm{sep}}(x)}(dy),
\qquad
Q_{1,P_{\mathrm{sep}}}(dy\mid x)
=\delta_{m_{0,\mathrm{sep}}(x)+\tau_{\mathrm{sep}}(x)}(dy),
\]
\[
\int y\,Q_{0,P_{\mathrm{sep}}}(dy\mid x)=m_{0,\mathrm{sep}}(x),
\qquad
\int y\,Q_{1,P_{\mathrm{sep}}}(dy\mid x)
=m_{0,\mathrm{sep}}(x)+\tau_{\mathrm{sep}}(x).
\]
The sine and cosine bounds give
\[
\frac14\le e_{\mathrm{sep}}(x)\le\frac34,
\qquad
|m_{0,\mathrm{sep}}(x)|=|\tau_{\mathrm{sep}}(x)|\le\frac18.
\]
Their increment bounds are
\[
\begin{aligned}
|e_{\mathrm{sep}}(x)-e_{\mathrm{sep}}(z)|
&\le\frac{\pi}{5}|x-z|
\le20|x-z|,\\
|m_{0,\mathrm{sep}}(x)-m_{0,\mathrm{sep}}(z)|
=|\tau_{\mathrm{sep}}(x)-\tau_{\mathrm{sep}}(z)|
&\le\frac{\pi}{4}|x-z|
\le20|x-z|.
\end{aligned}
\]
For every exponent \(\sigma_{\mathrm H}\in(0,1]\),
\[
|x-z|\le |x-z|^{\sigma_{\mathrm H}}
\qquad(x,z\in[0,1]).
\]
Consequently these continuous functions satisfy the supremum and increment bounds in
\cref{obj:ass:propensity-smoothness,obj:ass:baseline-smoothness,obj:ass:effect-smoothness}
at their respective exponents.

For the moment restriction, the baseline and effect caps imply
\[
|m_{0,\mathrm{sep}}(x)|\le1,
\qquad
|m_{0,\mathrm{sep}}(x)+\tau_{\mathrm{sep}}(x)|
\le |m_{0,\mathrm{sep}}(x)|+|\tau_{\mathrm{sep}}(x)|
\le1.
\]
Since \(p>0\), both deterministic arm kernels therefore satisfy
\[
\int |y|^p\,Q_{a_{\mathrm{arm}},P_{\mathrm{sep}}}(dy\mid x)
\le1\le10
\qquad
\bigl(a_{\mathrm{arm}}\in\{0,1\},\ x\in[0,1]\bigr).
\]
This verifies every restriction in \cref{obj:def:model}, and hence
\[
P_{\mathrm{sep}}\in\mathcal M_v.
\]
% lean: CausalSmith.Stat.FinitepHomogeneityDensegamma.separatedWitness_inModel

\item\label{proof:causal-nonempty:distance-bounds} For each \(P\in\mathcal M_v\), define its average contrast by
\[
\bar\tau_P=\int_0^1\tau_P(x)\,d\lambda(x).
\]
By \cref{obj:ass:effect-cap},
\[
|\bar\tau_P|
\le\int_0^1|\tau_P(x)|\,d\lambda(x)
\le\frac12,
\qquad
|\tau_P(x)-\bar\tau_P|\le1.
\]
The bounded functions below are integrable. Expanding the square and using
\(\lambda([0,1])=1\) gives
\[
\begin{aligned}
\int_0^1(\tau_P(x)-\bar\tau_P)^2\,d\lambda(x)
&=\int_0^1\tau_P(x)^2\,d\lambda(x)
 -2\bar\tau_P\int_0^1\tau_P(x)\,d\lambda(x)
 +\bar\tau_P^2\\
&=\int_0^1\tau_P(x)^2\,d\lambda(x)-\bar\tau_P^2.
\end{aligned}
\]
Moreover,
\[
\int_0^1\tau_P(x)^2\,d\lambda(x)
\le\int_0^1\frac14\,d\lambda(x)=\frac14.
\]
Thus
\[
0\le d(P)\le\frac12
\qquad(P\in\mathcal M_v).
\]
The distance set is therefore bounded above. It is nonempty by the witness from the preceding step, so its supremum satisfies
\[
D_v\le\frac12.
\]

For the witness, substitution in the cosine integrals gives
\[
\bar\tau_{P_{\mathrm{sep}}}
=\frac{1}{16\pi}\int_0^{2\pi}\cos\vartheta\,d\vartheta
=\frac{1}{16\pi}[\sin\vartheta]_0^{2\pi}
=0
\]
and
\[
\begin{aligned}
\int_0^1\tau_{\mathrm{sep}}(x)^2\,d\lambda(x)
&=\frac{1}{128\pi}
  \int_0^{2\pi}\cos^2\vartheta\,d\vartheta\\
&=\frac{1}{128\pi}
 \left[\frac{\vartheta}{2}
       +\frac{\sin\vartheta\cos\vartheta}{2}\right]_0^{2\pi}
=\frac1{128}.
\end{aligned}
\]
Since
\[
d_0>0,\qquad d_0^2=\frac1{128},
\]
taking the nonnegative square root yields
\[
d(P_{\mathrm{sep}})=d_0.
\]
Membership of this distance in the bounded distance set then gives
\[
d_0=d(P_{\mathrm{sep}})\le D_v\le\frac12.
\]
% lean: CausalSmith.Stat.FinitepHomogeneityDensegamma.model_distance_upper
% lean: CausalSmith.Stat.FinitepHomogeneityDensegamma.separatedWitness_distance
% lean: CausalSmith.Stat.FinitepHomogeneityDensegamma.model_distance_lower

\item Define the constant witness by
\[
e_{\mathrm{const}}=e_{\mathrm{sep}},
\qquad
m_{0,\mathrm{const}}=m_{0,\mathrm{sep}},
\qquad
\tau_{\mathrm{const}}(x)=\frac18
\quad(x\in[0,1]),
\]
\[
\begin{aligned}
P_{\mathrm{const}}(dx,da,dy)
=\lambda(dx)\bigl[
 &(1-e_{\mathrm{const}}(x))\,
 \delta_0(da)\delta_{m_{0,\mathrm{const}}(x)}(dy)\\
 &+e_{\mathrm{const}}(x)\,
 \delta_1(da)\delta_{m_{0,\mathrm{const}}(x)+1/8}(dy)
\bigr].
\end{aligned}
\]
The constant contrast satisfies
\[
|\tau_{\mathrm{const}}(x)|=\frac18\le\frac12,
\qquad
|\tau_{\mathrm{const}}(x)-\tau_{\mathrm{const}}(z)|
=0\le20|x-z|^\gamma.
\]
The preceding deterministic construction again gives uniform design, the displayed primitives, and the required arm means. The propensity and baseline restrictions have already been verified. The same moment argument applies because
\[
|m_{0,\mathrm{const}}(x)|\le\frac12,
\qquad
|m_{0,\mathrm{const}}(x)+\tau_{\mathrm{const}}(x)|\le1.
\]
Thus \(P_{\mathrm{const}}\in\mathcal M_v\), and its contrast is everywhere the admissible constant \(1/8\). By \cref{obj:def:null},
\[
P_{\mathrm{const}}\in H_0(v).
\]
% lean: CausalSmith.Stat.FinitepHomogeneityDensegamma.constantWitness_inNull
% lean: CausalSmith.Stat.FinitepHomogeneityDensegamma.deterministicLaw_primitives

\item Fix \(P\in\mathcal M_v\), and consider any other observational representation of the same probability law. Denote its continuous primitives and arm kernels by
\[
\widetilde e_P,\widetilde m_{0,P},\widetilde\tau_P\in C([0,1]),
\qquad
\widetilde m_{1,P}
=\widetilde m_{0,P}+\widetilde\tau_P,
\]
\[
\widetilde Q_{a_{\mathrm{arm}},P}(\,\cdot\mid x)
\in\mathcal P(\mathbb R)
\qquad
\bigl(a_{\mathrm{arm}}\in\{0,1\},\ x\in[0,1]\bigr),
\]
where \(\mathcal P(\mathbb R)\) denotes the Borel probability measures on \(\mathbb R\).
Define the two conditional record kernels by
\[
\begin{aligned}
\mathsf K_P(x,da,dy)
&=(1-e_P(x))\delta_0(da)Q_{0,P}(dy\mid x)
  +e_P(x)\delta_1(da)Q_{1,P}(dy\mid x),\\
\widetilde{\mathsf K}_P(x,da,dy)
&=(1-\widetilde e_P(x))\delta_0(da)\widetilde Q_{0,P}(dy\mid x)
  +\widetilde e_P(x)\delta_1(da)\widetilde Q_{1,P}(dy\mid x).
\end{aligned}
\]
The second representation has covariate marginal \(\lambda\), because its probability law is \(P\). Their observational factorizations therefore give
\[
P(dx,da,dy)
=\lambda(dx)\mathsf K_P(x,da,dy)
=\lambda(dx)\widetilde{\mathsf K}_P(x,da,dy).
\]
Uniqueness of probability kernels in this factorization implies
\[
\widetilde{\mathsf K}_P(x,\cdot)
=\mathsf K_P(x,\cdot)
\qquad\text{for \(\lambda\)-almost every }x.
\]
Evaluating the equality on \(\{1\}\times\mathbb R\) yields
\[
\widetilde e_P(x)=e_P(x)
\qquad\text{for \(\lambda\)-almost every }x.
\]
Continuous functions agreeing \(\lambda\)-almost everywhere agree throughout \([0,1]\): a nonzero difference at any point would, by continuity, remain nonzero on a relative neighborhood of positive Lebesgue measure. Hence
\[
\widetilde e_P=e_P
\qquad\text{pointwise on }[0,1].
\]

For every Borel \(\mathsf E\subseteq\mathbb R\), evaluation on the two arm slices now gives, at almost every \(x\),
\[
\begin{aligned}
(1-e_P(x))\widetilde Q_{0,P}(\mathsf E\mid x)
&=(1-e_P(x))Q_{0,P}(\mathsf E\mid x),\\
e_P(x)\widetilde Q_{1,P}(\mathsf E\mid x)
&=e_P(x)Q_{1,P}(\mathsf E\mid x).
\end{aligned}
\]
Both weights are strictly positive by \cref{obj:ass:overlap}. Cancelling them proves
\[
\widetilde Q_{0,P}(\,\cdot\mid x)=Q_{0,P}(\,\cdot\mid x),
\qquad
\widetilde Q_{1,P}(\,\cdot\mid x)=Q_{1,P}(\,\cdot\mid x)
\qquad\text{for \(\lambda\)-almost every }x.
\]
The arm-mean version identities consequently imply
\[
\begin{aligned}
\widetilde m_{0,P}(x)
&=\int y\,\widetilde Q_{0,P}(dy\mid x)
 =\int y\,Q_{0,P}(dy\mid x)
 =m_{0,P}(x),\\
\widetilde m_{1,P}(x)
&=\int y\,\widetilde Q_{1,P}(dy\mid x)
 =\int y\,Q_{1,P}(dy\mid x)
 =m_{1,P}(x)
\end{aligned}
\qquad\text{for \(\lambda\)-almost every }x.
\]
Continuity upgrades both equalities to pointwise equalities. Subtraction then gives
\[
\widetilde\tau_P
=\widetilde m_{1,P}-\widetilde m_{0,P}
=m_{1,P}-m_{0,P}
=\tau_P
\qquad\text{pointwise on }[0,1].
\]
% lean: CausalSmith.Stat.FinitepHomogeneityDensegamma.observed_continuous_versions_unique

\item The smoothness restrictions in
\cref{obj:ass:baseline-smoothness,obj:ass:effect-smoothness}, together with
\(m_{1,P}=m_{0,P}+\tau_P\), give
\[
\begin{aligned}
|m_{1,P}(x)-m_{1,P}(z)|
&=|(m_{0,P}(x)-m_{0,P}(z))
   +(\tau_P(x)-\tau_P(z))|\\
&\le |m_{0,P}(x)-m_{0,P}(z)|
    +|\tau_P(x)-\tau_P(z)|\\
&\le20|x-z|^\beta+20|x-z|^\gamma\\
&\le40|x-z|^{\min\{\beta,\gamma\}}.
\end{aligned}
\]
The last inequality uses \(0\le|x-z|\le1\) and
\(\min\{\beta,\gamma\}>0\). It includes \(x=z\), when all displayed increments and powers vanish.
% lean: CausalSmith.Stat.FinitepHomogeneityDensegamma.treated_mean_holder

\item In the completion of \cref{obj:def:causal-completion}, define
\[
Y_{\mathrm{obs}}=(1-A)Y(0)+AY(1).
\]
For any nonnegative Borel function
\[
f_{\mathrm{obs}}:[0,1]\times\{0,1\}\times\mathbb R
\longrightarrow[0,\infty],
\]
integration over the completion gives
\[
\begin{aligned}
&\int f_{\mathrm{obs}}(X,A,Y_{\mathrm{obs}})\,d\mathbb Q_P\\
&=\int_0^1
 \int_{\mathbb R}\int_{\mathbb R}
 \bigl[
 (1-e_P(x))f_{\mathrm{obs}}(x,0,y_0)
 +e_P(x)f_{\mathrm{obs}}(x,1,y_1)
 \bigr]\,
 Q_{1,P}(dy_1\mid x)Q_{0,P}(dy_0\mid x)\,d\lambda(x)\\
&=\int_0^1\left[
 (1-e_P(x))\int_{\mathbb R}f_{\mathrm{obs}}(x,0,y)\,Q_{0,P}(dy\mid x)
 +e_P(x)\int_{\mathbb R}f_{\mathrm{obs}}(x,1,y)\,Q_{1,P}(dy\mid x)
 \right]d\lambda(x)\\
&=\int f_{\mathrm{obs}}\,dP.
\end{aligned}
\]
Indeed, in each arm the unused potential-outcome kernel integrates to one; the last equality is the observational factorization with uniform covariate marginal. Equality for all such functions proves
\[
\mathcal L_{\mathbb Q_P}(X,A,Y_{\mathrm{obs}})=P.
\]
% lean: CausalSmith.Stat.FinitepHomogeneityDensegamma.completion_observed_marginal

\item\label{proof:causal-nonempty:conditional-law} Integrating out the probability kernels first gives
\[
\mathcal L_{\mathbb Q_P}(X)=\lambda.
\]
The conditional law of the remaining coordinates is the defining product kernel:
\[
\begin{aligned}
&\mathcal L_{\mathbb Q_P}(A,Y(0),Y(1)\mid X=x)\\
&\quad=
\bigl((1-e_P(x))\delta_0+e_P(x)\delta_1\bigr)
\otimes Q_{0,P}(\,\cdot\mid x)
\otimes Q_{1,P}(\,\cdot\mid x)
\qquad\text{for \(\lambda\)-almost every }x.
\end{aligned}
\]
Projecting this identity onto treatment and onto the potential-outcome pair yields
\[
\begin{aligned}
\mathcal L_{\mathbb Q_P}(A\mid X=x)
&=(1-e_P(x))\delta_0+e_P(x)\delta_1,\\
\mathcal L_{\mathbb Q_P}(Y(0),Y(1)\mid X=x)
&=Q_{0,P}(\,\cdot\mid x)\otimes Q_{1,P}(\,\cdot\mid x).
\end{aligned}
\]
Thus the joint conditional law equals the product of these two conditional marginals. The conditional-independence criterion gives
\[
(Y(0),Y(1))\perp A\mid X.
\]
% lean: CausalSmith.Stat.FinitepHomogeneityDensegamma.completion_conditional_independence

\item For \(y\in\mathbb R\), the elementary bound needed for integrability is
\[
|y|\le1+|y|^p.
\]
If \(|y|\le1\), this follows from \(|y|^p\ge0\). If \(|y|>1\), the condition \(p>1\) gives
\[
|y|\le|y|^p\le1+|y|^p.
\]
Hence \cref{obj:ass:raw-moment} implies, for each
\(a_{\mathrm{arm}}\in\{0,1\}\),
\[
\int_{\mathbb R}|y|\,Q_{a_{\mathrm{arm}},P}(dy\mid x)
\le1+\int_{\mathbb R}|y|^p\,Q_{a_{\mathrm{arm}},P}(dy\mid x)
\le11
\qquad\text{for \(\lambda\)-almost every }x.
\]
Integrating the assignment kernel and the unused potential-outcome kernel, both of total mass one, gives
\[
\begin{aligned}
\int |Y(a_{\mathrm{arm}})|\,d\mathbb Q_P
&=\int_0^1\int_{\mathbb R}
 |y|\,Q_{a_{\mathrm{arm}},P}(dy\mid x)\,d\lambda(x)\\
&\le\int_0^1 11\,d\lambda(x)=11<\infty.
\end{aligned}
\]
This proves integrability for \(Y(0)\) and \(Y(1)\).
% lean: CausalSmith.Stat.FinitepHomogeneityDensegamma.completion_potential_integrable

\item For each arm, integration against the joint conditional law from \cref{proof:causal-nonempty:conditional-law} reduces to the corresponding arm mean, since the other two factors are probability measures. The conditional-expectation formula for an integrable function therefore gives
\[
\mathbb E_{\mathbb Q_P}[Y(a_{\mathrm{arm}})\mid X]
=\int_{\mathbb R}y\,Q_{a_{\mathrm{arm}},P}(dy\mid X)
\qquad\mathbb Q_P\text{-almost surely},
\quad a_{\mathrm{arm}}\in\{0,1\}.
\]
Because \(X\) has law \(\lambda\), the observational arm-mean version identities hold also at \(X\), almost surely under \(\mathbb Q_P\):
\[
\int y\,Q_{0,P}(dy\mid X)=m_{0,P}(X),
\qquad
\int y\,Q_{1,P}(dy\mid X)=m_{0,P}(X)+\tau_P(X).
\]
Linearity of conditional expectation, justified by the integrability established above, now yields
\[
\begin{aligned}
\mathbb E_{\mathbb Q_P}[Y(1)-Y(0)\mid X]
&=\mathbb E_{\mathbb Q_P}[Y(1)\mid X]
 -\mathbb E_{\mathbb Q_P}[Y(0)\mid X]\\
&=(m_{0,P}(X)+\tau_P(X))-m_{0,P}(X)
=\tau_P(X)
\end{aligned}
\qquad\mathbb Q_P\text{-almost surely}.
\]
% lean: CausalSmith.Stat.FinitepHomogeneityDensegamma.completion_conditional_effect

\item Finally, suppose \(d(P)=0\). The defining integral is nonnegative, so
\[
\int_0^1(\tau_P(x)-\bar\tau_P)^2\,d\lambda(x)=0.
\]
Its integrand is nonnegative and bounded by one, as established in \cref{proof:causal-nonempty:distance-bounds}. The zero-integral criterion therefore gives
\[
(\tau_P(x)-\bar\tau_P)^2=0
\quad\text{and hence}\quad
\tau_P(x)=\bar\tau_P
\qquad\text{for \(\lambda\)-almost every }x.
\]
Continuity and the full support of \(\lambda\) upgrade this to
\[
\tau_P(x)=\bar\tau_P
\qquad(x\in[0,1]).
\]
The average bound in \cref{proof:causal-nonempty:distance-bounds} gives
\[
-\frac12\le\bar\tau_P\le\frac12.
\]
Together with \(P\in\mathcal M_v\), this is exactly the membership condition in \cref{obj:def:null}.

Conversely, if \(P\in H_0(v)\), that definition supplies a constant
\[
c_{\mathrm{null}}\in[-1/2,1/2],
\qquad
\tau_P(x)=c_{\mathrm{null}}
\quad(x\in[0,1]).
\]
Since \(\lambda\) is a probability measure,
\[
\bar\tau_P=c_{\mathrm{null}},
\qquad
d(P)
=\left\{\int_0^1(c_{\mathrm{null}}-c_{\mathrm{null}})^2
\,d\lambda(x)\right\}^{1/2}
=0.
\]
Thus
\[
d(P)=0\quad\Longleftrightarrow\quad P\in H_0(v),
\]
completing the proof.
% lean: CausalSmith.Stat.FinitepHomogeneityDensegamma.hetDist_zero_iff_null
\end{enumerate}
\end{proof}
